\section{Results and Discussion}
\label{sec:results_and_discussion}

\subsection{Global Performance}
\label{sec:global_performance}
\label{sec:global_performance_analysis}
\label{sec:global_classification_performance}

The empirical evaluation of the RF classifier across seven feature configuration schemes on the held-out test set under the Strict Subject-Disjoint protocol is summarized in Table~\ref{tab:VII}. The experimental results demonstrate that the Tri-Domain configuration achieves the highest performance with an Accuracy of 85.00\%, Precision of 85.06\%, Recall of 85.00\%, and F1-Score of 84.93\%. This achievement outperforms the dual-domain Emotion $\oplus$ Face scheme (Accuracy of 84.86\% and F1-Score of 84.79\%) and exceeds the single-domain Face baseline of 84.03\%. The Age domain records the lowest accuracy (71.81\%), reflecting the challenges of partitioning high-dimensional continuous features via orthogonal random tree splits. The majority of configurations select PCA 0.75 dimensionality reduction, confirming that increasing feature dimensions without representation selection increases the variance of random feature selection at tree nodes~\cite{ref48}.

\begin{table*}[htbp]
\caption{Performance Benchmark of Random Forest across Seven Feature Configurations.}
\label{tab:7}
\label{tab:VII}
\centering
\footnotesize
\setlength{\tabcolsep}{2.5pt}
\resizebox{\textwidth}{!}{%
\begin{tabular}{lccccccc}
\toprule
\textbf{Feature Configuration} & \textbf{CV Score (\%)} & \textbf{Accuracy (\%)} & \textbf{95\% CI} & \textbf{Precision (\%)} & \textbf{Recall (\%)} & \textbf{F1-Score (\%)} & \textbf{Best Hyperparameters} \\
\midrule
Face & 85.97 & 84.03 & [82.42, 85.51] & 84.13 & 84.03 & 84.03 & \begin{tabular}[t]{@{}l@{}}n\_est=200, depth=20, max\_feat=log2, \\ min\_split=5, min\_leaf=1, \\ pca=PCA(0.75), scaler=None\end{tabular} \\
Emotion & 77.41 & 76.11 & [74.27, 77.86] & 76.02 & 76.11 & 76.03 & \begin{tabular}[t]{@{}l@{}}n\_est=200, depth=None, max\_feat=sqrt, \\ min\_split=2, min\_leaf=2, \\ pca=PCA(0.75), scaler=None\end{tabular} \\
Age & 70.36 & 71.81 & [69.87, 73.66] & 71.76 & 71.81 & 71.70 & \begin{tabular}[t]{@{}l@{}}n\_est=200, depth=None, max\_feat=log2, \\ min\_split=5, min\_leaf=2, \\ pca=PCA(0.75), scaler=None\end{tabular} \\
Emotion $\oplus$ Face & 86.09 & 84.86 & [83.29, 86.31] & 84.93 & 84.86 & 84.79 & \begin{tabular}[t]{@{}l@{}}n\_est=200, depth=None, max\_feat=sqrt, \\ min\_split=2, min\_leaf=1, \\ pca=PCA(0.75), scaler=MinMaxScaler\end{tabular} \\
Face $\oplus$ Age & 85.76 & 84.86 & [83.29, 86.31] & 84.88 & 84.86 & 84.78 & \begin{tabular}[t]{@{}l@{}}n\_est=200, depth=None, max\_feat=sqrt, \\ min\_split=2, min\_leaf=1, \\ pca=None, scaler=None\end{tabular} \\
Emotion $\oplus$ Age & 77.03 & 77.31 & [75.50, 79.03] & 77.34 & 77.31 & 77.28 & \begin{tabular}[t]{@{}l@{}}n\_est=200, depth=20, max\_feat=log2, \\ min\_split=2, min\_leaf=2, \\ pca=PCA(0.75), scaler=None\end{tabular} \\
\textbf{Face $\oplus$ Emotion $\oplus$ Age} & \textbf{85.81} & \textbf{85.00} & \textbf{[83.43, 86.44]} & \textbf{85.06} & \textbf{85.00} & \textbf{84.93} & \begin{tabular}[t]{@{}l@{}}\textbf{n\_est=200, depth=30, max\_feat=sqrt,} \\ \textbf{min\_split=2, min\_leaf=1,} \\ \textbf{pca=None, scaler=MinMaxScaler}\end{tabular} \\
\bottomrule
\end{tabular}%
}
\end{table*}

The performance evaluation of the GNB classifier across seven feature configurations on the held-out test set is presented in Table~\ref{tab:VIII}. Unlike the other classifiers, the peak performance of GNB is achieved by the dual-domain Emotion $\oplus$ Face configuration with an Accuracy of 84.31\%, Precision of 84.37\%, Recall of 84.31\%, and F1-Score of 84.30\%. In the Tri-Domain configuration, classification performance declines to an Accuracy of 83.61\% and an F1-Score of 83.64\%. This degradation is a direct consequence of violating the class-conditional feature independence assumption underlying the multivariate Gaussian distribution. When the dimensionality is expanded to 2,304 dimensions through the concatenation of all three domains, latent inter-feature correlations increase sharply, thereby distorting univariate probability estimations. Meanwhile, the single-domain Face category records an Accuracy of 82.55\%, surpassing Emotion (72.13\%) and Age (67.08\%). All configurations uniformly select PCA 0.75 dimensionality reduction.

\begin{table*}[htbp]
\caption{Performance Benchmark of Gaussian Naive Bayes across Seven Feature Configurations.}
\label{tab:8}
\label{tab:VIII}
\centering
\footnotesize
\setlength{\tabcolsep}{2.5pt}
\resizebox{\textwidth}{!}{%
\begin{tabular}{lccccccc}
\toprule
\textbf{Feature Configuration} & \textbf{CV Score (\%)} & \textbf{Accuracy (\%)} & \textbf{95\% CI} & \textbf{Precision (\%)} & \textbf{Recall (\%)} & \textbf{F1-Score (\%)} & \textbf{Best Hyperparameters} \\
\midrule
Face & 83.43 & 82.55 & [80.89, 84.09] & 82.68 & 82.55 & 82.50 & \begin{tabular}[t]{@{}l@{}}var\_smoothing=4.1246e-02, \\ pca=PCA(0.75), scaler=MinMaxScaler\end{tabular} \\
Emotion & 71.42 & 72.13 & [70.20, 73.98] & 72.51 & 72.13 & 71.99 & \begin{tabular}[t]{@{}l@{}}var\_smoothing=5.8780e-03, \\ pca=PCA(0.75), scaler=MinMaxScaler\end{tabular} \\
Age & 68.18 & 67.08 & [65.07, 69.03] & 66.91 & 67.08 & 66.86 & \begin{tabular}[t]{@{}l@{}}var\_smoothing=8.3768e-04, \\ pca=PCA(0.75), scaler=MinMaxScaler\end{tabular} \\
\textbf{Emotion $\oplus$ Face} & \textbf{84.65} & \textbf{84.31} & \textbf{[82.71, 85.78]} & \textbf{84.37} & \textbf{84.31} & \textbf{84.30} & \begin{tabular}[t]{@{}l@{}}\textbf{var\_smoothing=1.1253e-02,} \\ \textbf{pca=PCA(0.75), scaler=MinMaxScaler}\end{tabular} \\
Face $\oplus$ Age & 83.48 & 83.06 & [81.42, 84.58] & 83.47 & 83.06 & 83.10 & \begin{tabular}[t]{@{}l@{}}var\_smoothing=1.1253e-02, \\ pca=PCA(0.75), scaler=MinMaxScaler\end{tabular} \\
Emotion $\oplus$ Age & 73.07 & 75.88 & [74.03, 77.64] & 75.97 & 75.88 & 75.78 & \begin{tabular}[t]{@{}l@{}}var\_smoothing=1.1253e-02, \\ pca=PCA(0.75), scaler=MinMaxScaler\end{tabular} \\
Face $\oplus$ Emotion $\oplus$ Age & 84.41 & 83.61 & [81.99, 85.11] & 83.79 & 83.61 & 83.64 & \begin{tabular}[t]{@{}l@{}}var\_smoothing=1.1253e-02, \\ pca=PCA(0.75), scaler=None\end{tabular} \\
\bottomrule
\end{tabular}%
}
\end{table*}

The empirical benchmark results of the multinomial LR classifier across seven feature configurations on the held-out test set are presented in Table~\ref{tab:IX}. As the best-performing model, Tri-Domain LR (Face $\oplus$ Emotion $\oplus$ Age) achieves an Accuracy of 88.66\%, Precision of 88.75\%, Recall of 88.66\%, and F1-Score of 88.68\%. This achievement consistently surpasses the single-domain baselines of Face (Accuracy of 86.39\%), Emotion (85.23\%), and Age (81.48\%), while also outperforming the dual-domain Emotion $\oplus$ Face scheme (87.45\%). The success of LR demonstrates the strong linear complementarity across the spans of latent representations. Through an $L_2$-regularized Softmax activation function, LR effectively models high-dimensional multivariable feature covariates~\cite{ref49}. The linear weights of the model establish a stable multiclass separating hyperplane decision boundary via weighted regularization~\cite{ref50}, where all configurations retain their full dimensionality without PCA reduction.

\begin{table*}[htbp]
\caption{Performance Benchmark of Logistic Regression across Seven Feature Configurations.}
\label{tab:9}
\label{tab:IX}
\centering
\footnotesize
\setlength{\tabcolsep}{2.5pt}
\resizebox{\textwidth}{!}{%
\begin{tabular}{lccccccc}
\toprule
\textbf{Feature Configuration} & \textbf{CV Score (\%)} & \textbf{Accuracy (\%)} & \textbf{95\% CI} & \textbf{Precision (\%)} & \textbf{Recall (\%)} & \textbf{F1-Score (\%)} & \textbf{Best Hyperparameters} \\
\midrule
Face & 87.99 & 86.39 & [84.88, 87.77] & 86.46 & 86.39 & 86.40 & \begin{tabular}[t]{@{}l@{}}C=0.1, solver=saga, max\_iter=500, \\ pca=None, scaler=None\end{tabular} \\
Emotion & 85.61 & 85.23 & [83.67, 86.67] & 85.35 & 85.23 & 85.27 & \begin{tabular}[t]{@{}l@{}}C=0.1, solver=lbfgs, max\_iter=500, \\ pca=None, scaler=None\end{tabular} \\
Age & 83.23 & 81.48 & [79.79, 83.06] & 81.45 & 81.48 & 81.46 & \begin{tabular}[t]{@{}l@{}}C=0.1, solver=lbfgs, max\_iter=500, \\ pca=None, scaler=None\end{tabular} \\
Emotion $\oplus$ Face & 89.54 & 87.45 & [85.99, 88.78] & 87.59 & 87.45 & 87.49 & \begin{tabular}[t]{@{}l@{}}C=1, solver=saga, max\_iter=500, \\ pca=None, scaler=MinMaxScaler\end{tabular} \\
Face $\oplus$ Age & 89.12 & 87.41 & [85.94, 88.74] & 87.44 & 87.41 & 87.40 & \begin{tabular}[t]{@{}l@{}}C=0.1, solver=lbfgs, max\_iter=500, \\ pca=None, scaler=MinMaxScaler\end{tabular} \\
Emotion $\oplus$ Age & 86.96 & 86.67 & [85.17, 88.04] & 86.68 & 86.67 & 86.67 & \begin{tabular}[t]{@{}l@{}}C=0.1, solver=saga, max\_iter=500, \\ pca=None, scaler=MinMaxScaler\end{tabular} \\
\textbf{Face $\oplus$ Emotion $\oplus$ Age} & \textbf{90.00} & \textbf{88.66} & \textbf{[87.25, 89.93]} & \textbf{88.75} & \textbf{88.66} & \textbf{88.68} & \begin{tabular}[t]{@{}l@{}}\textbf{C=0.1, solver=saga, max\_iter=500,} \\ \textbf{pca=None, scaler=MinMaxScaler}\end{tabular} \\
\bottomrule
\end{tabular}%
}
\end{table*}

The performance progression of the SVM classifier across all feature configurations is presented in Table~\ref{tab:X}. As the second-best model, Tri-Domain SVM (Face $\oplus$ Emotion $\oplus$ Age) achieves an Accuracy of 88.61\%, Precision of 88.71\%, Recall of 88.61\%, and F1-Score of 88.63\% using a standard RBF kernel formulation~\cite{ref51} with a penalty parameter of $C = 10$. The evaluation reveals a significant performance surge from the single-domain Face baseline (Accuracy of 86.16\%) to the dual-domain Emotion $\oplus$ Face configuration (Accuracy of 88.43\%), representing a net gain of 2.27 percentage points (pp). However, incorporating age features into the Tri-Domain scheme yields only a marginal improvement of 0.18~pp. This phenomenon indicates asymptotic representational saturation in the non-linear Hilbert space mapping of the RBF kernel, where the joint biometric and emotion manifold has already extracted nearly all primary discriminative variance across intersectional subgroups.

\begin{table*}[htbp]
\caption{Performance Benchmark of Support Vector Machine across Seven Feature Configurations.}
\label{tab:10}
\label{tab:X}
\centering
\footnotesize
\setlength{\tabcolsep}{2.5pt}
\resizebox{\textwidth}{!}{%
\begin{tabular}{lccccccc}
\toprule
\textbf{Feature Configuration} & \textbf{CV Score (\%)} & \textbf{Accuracy (\%)} & \textbf{95\% CI} & \textbf{Precision (\%)} & \textbf{Recall (\%)} & \textbf{F1-Score (\%)} & \textbf{Best Hyperparameters} \\
\midrule
Face & 88.39 & 86.16 & [84.64, 87.55] & 86.24 & 86.16 & 86.19 & \begin{tabular}[t]{@{}l@{}}C=0.01, poly, $\gamma$=scale, deg=3, \\ pca=None, scaler=MinMaxScaler\end{tabular} \\
Emotion & 85.90 & 85.79 & [84.25, 87.20] & 85.84 & 85.79 & 85.80 & \begin{tabular}[t]{@{}l@{}}C=10, rbf, $\gamma$=scale, \\ pca=None, scaler=None\end{tabular} \\
Age & 83.67 & 82.55 & [80.89, 84.09] & 82.68 & 82.55 & 82.58 & \begin{tabular}[t]{@{}l@{}}C=10, rbf, $\gamma$=scale, \\ pca=None, scaler=None\end{tabular} \\
Emotion $\oplus$ Face & 89.39 & 88.43 & [87.01, 89.71] & 88.50 & 88.43 & 88.44 & \begin{tabular}[t]{@{}l@{}}C=0.1, poly, $\gamma$=scale, deg=2, \\ pca=None, scaler=MinMaxScaler\end{tabular} \\
Face $\oplus$ Age & 89.24 & 87.55 & [86.09, 88.87] & 87.66 & 87.55 & 87.57 & \begin{tabular}[t]{@{}l@{}}C=10, rbf, $\gamma$=scale, \\ pca=None, scaler=None\end{tabular} \\
Emotion $\oplus$ Age & 87.12 & 86.25 & [84.73, 87.64] & 86.26 & 86.25 & 86.24 & \begin{tabular}[t]{@{}l@{}}C=10, rbf, $\gamma$=scale, \\ pca=None, scaler=None\end{tabular} \\
\textbf{Face $\oplus$ Emotion $\oplus$ Age} & \textbf{89.86} & \textbf{88.61} & \textbf{[87.20, 89.88]} & \textbf{88.71} & \textbf{88.61} & \textbf{88.63} & \begin{tabular}[t]{@{}l@{}}\textbf{C=10, rbf, $\gamma$=scale,} \\ \textbf{pca=None, scaler=MinMaxScaler}\end{tabular} \\
\bottomrule
\end{tabular}%
}
\end{table*}

A comparative synthesis across models in Table~\ref{tab:VII}, Table~\ref{tab:VIII}, Table~\ref{tab:IX}, and Table~\ref{tab:X} confirms that the effectiveness of multi-domain representation fusion strongly depends on the classifier architecture. The linear LR and non-linear SVM models lead the global benchmark with competitive peak performance in the Tri-Domain configuration (Accuracy of 88.66\% vs. 88.61\%), differing marginally by only a single test image out of the 2,160 total samples. In contrast, RF and GNB exhibit distinctly different decision boundary dynamics. RF achieves 85.00\% accuracy in the Tri-Domain configuration but is constrained by the complexity of high-dimensional random tree partitioning, whereas GNB reaches its optimal point in the dual-domain scheme (84.31\%) due to the constraints of the covariance independence assumption. These findings demonstrate that multi-domain latent feature fusion is most effectively mapped by classifiers with global decision boundary estimation.

\begin{table*}[htbp]
\caption{Vision Transformer Feature Extraction Runtime Profile.}
\label{tab:11}
\label{tab:XI}
\centering
\footnotesize
\setlength{\tabcolsep}{6pt}
\begin{tabular}{lcc}
\toprule
\textbf{Feature Domain} & \textbf{Average Extraction Time per Image (ms)} & \textbf{Total Extraction Time (s)} \\
\midrule
Biometric Face & 39.74 & 429.22 (7m 09s) \\
Facial Emotion & 39.65 & 428.24 (7m 08s) \\
Facial Age & 45.11 & 487.17 (8m 07s) \\
\midrule
\textbf{Tri-Domain Fusion} & \textbf{124.50} & \textbf{1,344.64 (22m 25s)} \\
\bottomrule
\end{tabular}
\end{table*}

The computational efficiency profile is evaluated across two measurable aspects in Table~\ref{tab:XI} and Table~\ref{tab:XII}. GPU-accelerated offline ViT feature extraction records a duration of 39.74~ms for Face, 39.65~ms for Emotion, and 45.11~ms for Age, yielding a total tri-domain duration of 124.50~ms per image (a total of 22.41 minutes for 10,800 images with a cache footprint of 92.4~MB, or 30.8~MB per domain, alongside a test array RAM consumption of 18.98~MB). In 5-fold cross-validation training, GNB requires 25.49 minutes for dual-domain, followed by LR at 54.78 minutes, RF at 110.47 minutes, and SVM reaching 511.10 minutes, confirming that LR training is 9.3 times faster than SVM. During inference testing, the downstream LR classification head demonstrates superiority with a latency of only 0.0046~ms per sample, 1,027 times faster than SVM (4.6785~ms). In terms of memory consumption and model storage, the held-out test array consumes 18.98~MB of RAM for the tri-domain configuration and 12.66~MB for the dual-domain configuration, whereas the serialized LR model file requires only 0.11~MB, which is 330 times more compact than SVM (36.60~MB), establishing LR as the most optimal downstream architecture for practical deployment.

\begin{table*}[htbp]
\caption{Downstream Classifier Training and Testing Computational Profile.}
\label{tab:12}
\label{tab:XII}
\centering
\footnotesize
\setlength{\tabcolsep}{3.5pt}
\resizebox{\textwidth}{!}{%
\begin{tabular}{llcccc}
\toprule
\textbf{Classifier} & \textbf{Feature Scheme} & \textbf{5-Fold CV Training Time (s)} & \textbf{Inference Latency (ms/sample)} & \textbf{RAM Consumption (MB)} & \textbf{Model Size (MB)} \\
\midrule
\textbf{Logistic Regression} & \textbf{Tri-Domain} & \textbf{3,287.07 (54.78 min)} & \textbf{0.0046} & \textbf{18.98} & \textbf{0.11} \\
Support Vector Machine & Tri-Domain & 30,666.02 (511.10 min) & 4.6785 & 18.98 & 36.60 \\
Random Forest & Tri-Domain & 6,628.03 (110.47 min) & 0.0334 & 18.98 & 6.19 \\
Gaussian Naive Bayes & Dual-Domain (Emotion $\oplus$ Face) & 1,529.50 (25.49 min) & 0.0088 & 12.66 & 0.78 \\
\bottomrule
\end{tabular}%
}
\end{table*}
